% =============================================================================
%  lec19.tex —— 第 19 讲 笛卡尔纤维化（底本 PDF 第 146–152 页）
% =============================================================================
\lecture{19}{笛卡尔纤维化}
\begin{center}\small\color{black!55}（底本 PDF 第 146–152 页）\end{center}

本讲我们定义 $\infty$-范畴之间的笛卡尔纤维化与余笛卡尔纤维化。

\section{局部笛卡尔箭头}

\nitem\label{L19-1.1} 非正式地说，\textbf{笛卡尔纤维化}是函子 $p\colon\cC\to\cD$，使得纤维 $\{\cC_d\}_{d\in\cD}$ 以同伦连贯的方式反变函子性地依赖于 $d$。对任意对象 $c'\in\cC_{d'}$ 与底范畴 $\cD$ 中的态射 $g\colon d\to d'$，相应的函子 $g^{\dagger}\colon\cC_{d'}\to\cC_d$ 把 $c'$ 送到一个对象 $c\in\cC_d$，并配以 $\cC$ 中一个在 $g$ 上方的态射 $f\colon c\to c'$。我们可以把上述讨论描绘成下图。
\begin{center}
\begin{tikzcd}[column sep=3em,row sep=2.8em]
c \arrow[r,dashed,"f"] & c' \arrow[d,"p"] & \in\cC\\
d \arrow[r,"g"] & d' & \in\cD
\end{tikzcd}
\end{center}
态射 $f$ 可由下述泛性质刻画。

\begin{definition}\label{L19-1.2}
设 $p\colon\cC\to\cD$ 是拟范畴之间的内纤维化。若对朴素\footnote{这指的是普通范畴 $\Set_\Delta$ 中的纤维积。注意 $\cC_{p(X)}$ 是拟范畴，因为 $p$ 是内纤维化。}纤维
\[
\cC_{p(X)}\coloneqq\cC\times_{\cD}\{p(X)\}
\]
中的每个对象 $X'$，$\Spc$ 中的下述交换方块都是笛卡尔的，则称 $\cC$ 中的态射 $f\colon X\to Y$ 是\textbf{局部 $p$-笛卡尔的}，或\textbf{在 $\cD$ 上笛卡尔}：
\begin{equation}\label{L19-1.1eq}
\begin{tikzcd}[column sep=4.6em,row sep=3.4em]
\operatorname{Maps}_{\cC_{p(X)}}(X',X) \arrow[r,"f\circ-"] \arrow[d] &
\operatorname{Maps}_{\cC}(X',Y) \arrow[d]\\
\{*\} \arrow[r,"p(f)"] & \operatorname{Maps}_{\cD}(p(X'),p(Y)).
\end{tikzcd}
\end{equation}
若对任意 $Y\in\cC$ 与 $\cD$ 中的态射 $g\colon x\to p(Y)$，都存在 $\cC$ 中在 $g$ 上方的一个局部 $p$-笛卡尔态射 $f\colon X\to Y$，则称 $p$ 是\textbf{局部笛卡尔纤维化}。

对偶地，若 $\cC^{\op}$ 中相应的态射是局部 $p$-笛卡尔的，则称 $f$ 是\textbf{局部 $p$-余笛卡尔的}。若 $p^{\op}$ 是局部笛卡尔纤维化，则称 $p$ 是\textbf{局部余笛卡尔纤维化}。
\end{definition}

\yizhu{“笛卡尔”在这里是 Lurie 对相对纤维化中一类态射的称谓（HTT 第 2.4 节），与“笛卡尔积”“笛卡尔闭范畴”等术语中的“笛卡尔”含义不同。译文按术语表的冻结译法，把 $p$-Cartesian 与 locally $p$-Cartesian 分别译作“$p$-笛卡尔”与“局部 $p$-笛卡尔”；后者弱于前者，见练习~\ref{L19-2.22}。\textbf{[来源：A]}}

\begin{remark}\label{L19-1.3}
非正式地说，$f\colon X\to Y$ 是局部 $p$-笛卡尔的，当且仅当下述数据等价：
\begin{itemize}[leftmargin=2.2em,itemsep=0.1em]
\item 纤维 $\cC_{p(X)}$ 中的态射 $X'\to X$；
\item $\cC$ 中在 $p(f)$ 上方的态射 $X'\to Y$。
\end{itemize}
\end{remark}

\begin{remark}\label{L19-1.4}
严格说来，$\Spc$ 中的交换方块 \eqref{L19-1.1eq} 只在上同伦意义下良定义。然而，“是笛卡尔的”这一性质在同伦下不变。等价地，我们可以说 $h\Spc$ 中相应的方块（它是良定义的）是同伦拉回。注意右侧的竖直映射可以实现为 Kan 复形之间的 Kan 纤维化（HTT.2.4.4.1）：
\[
\Hom^{R}_{\cC}(X',Y)\to\Hom^{R}_{\cD}(p(X'),p(Y)).
\]
\end{remark}

\begin{exercise}\label{L19-1.5}
若 $f$ 是同构，则 $f$ 是局部 $p$-笛卡尔的。
\end{exercise}

\begin{exercise}\label{L19-1.6}
若 $f$ 同伦于 $f'$，则 $f$ 是局部 $p$-笛卡尔的，当且仅当 $f'$ 是。
\end{exercise}

\begin{warning}\label{L19-1.7}
朴素纤维积 $\cC_{p(X)}\coloneqq\cC\times_{\cD}\{p(X)\}$ 未必是同伦纤维积或 $\infty$-范畴纤维积。
\end{warning}

\begin{exercise}\label{L19-1.8}
证明 $p\colon\mathrm{N}_{\bullet}([1])\to\mathrm{N}_{\bullet}(\mathrm{Isom})$ 是内纤维化，其朴素纤维都等价于 $\Delta^0$，而同伦纤维都等价于 $\Delta^1$。
\end{exercise}

\begin{exercise}\label{L19-1.9}
考虑内纤维化 $p\colon\mathrm{N}_{\bullet}([1])\to\mathrm{N}_{\bullet}(\mathrm{Isom})$ 与 $p'\colon\mathrm{N}_{\bullet}([1])\to\mathrm{N}_{\bullet}([0])$。证明：源中每条边都是局部 $p$-笛卡尔的，而只有退化边是局部 $p'$-笛卡尔的。
\end{exercise}

\begin{warning}\label{L19-1.10}
上述练习表明，定义~\ref{L19-1.2} 在范畴等价下并非不变。然而注意，若 $p\colon\cC\to\cD$ 是范畴纤维化，这不会发生，因为其朴素纤维与同伦纤维一致。
\end{warning}

\section{笛卡尔箭头}

\nitem\label{L19-2.1} 下述练习说明，局部 $p$-笛卡尔箭头可能对复合不封闭。因此我们需要一个更强的条件，才能使纤维具有函子性。

\begin{example}\label{L19-2.2}
考虑偏序集之间的下述映射
\[
[1]\times[1]\to[2],\qquad (0,0)\mapsto 0,\ (0,1)\mapsto 0,\ (1,0)\mapsto 1,\ (1,1)\mapsto 2.
\]
证明 $p\colon\mathrm{N}_{\bullet}([1]\times[1])\to\mathrm{N}_{\bullet}([2])$ 是局部笛卡尔纤维化，但局部 $p$-笛卡尔箭头对复合不封闭。
\end{example}

\begin{proposition}\label{L19-2.3}（HTT.2.4.2.7）
设 $p\colon\cC\to\cD$ 是拟范畴之间的局部笛卡尔纤维化，$f\colon X\to Y$ 是 $\cC$ 中的态射。则下述条件等价：
\begin{enumerate}[label=(\arabic*),leftmargin=2.4em,itemsep=0.1em]
\item 对任意 $X'\in\cC$，下述交换方块是笛卡尔的：
\begin{equation}\label{L19-2.1eq}
\begin{tikzcd}[column sep=4.6em,row sep=3.4em]
\operatorname{Maps}_{\cC}(X',X) \arrow[r,"f\circ-"] \arrow[d] &
\operatorname{Maps}_{\cC}(X',Y) \arrow[d]\\
\operatorname{Maps}_{\cD}(p(X'),p(X)) \arrow[r,"p(f)"] &
\operatorname{Maps}_{\cD}(p(X'),p(Y)).
\end{tikzcd}
\end{equation}
\item 对任意态射 $g\colon X'\to X$，$g$ 是局部 $p$-笛卡尔的，当且仅当 $f\circ g$ 是局部 $p$-笛卡尔的。
\item 对任意态射 $g\colon X'\to X$，若 $g$ 是局部 $p$-笛卡尔的，则 $f\circ g$ 是局部 $p$-笛卡尔的。
\end{enumerate}
\end{proposition}

\begin{remark}\label{L19-2.4}
非正式地说，(1) 意味着下述数据等价：
\begin{itemize}[leftmargin=2.2em,itemsep=0.1em]
\item $\cC$ 中的态射 $X'\to X$；
\item $\cC$ 中的态射 $X'\to Y$，以及 $p(X')\to p(Y)$ 经 $p(X)$ 的一个分解。
\end{itemize}
\end{remark}

\begin{definition}\label{L19-2.5}
设 $p\colon\cC\to\cD$ 是拟范畴之间的内纤维化，$f\colon X\to Y$ 是 $\cC$ 中的态射。若 $f$ 满足命题~\ref{L19-2.3} 中的条件 (1)，则称 $f$ 是\textbf{$p$-笛卡尔的}，或\textbf{在 $\cD$ 上笛卡尔}。若对任意 $Y\in\cC$ 与 $\cD$ 中的边 $g\colon x\to p(Y)$，都存在 $\cC$ 中在 $g$ 上方的一个 $p$-笛卡尔态射，则称 $p$ 是\textbf{笛卡尔纤维化}。

对偶地，若 $\cC^{\op}$ 中相应的箭头是 $p^{\op}$-笛卡尔的，则称 $f$ 是\textbf{$p$-余笛卡尔的}。若 $p^{\op}$ 是笛卡尔纤维化，则称 $p$ 是\textbf{余笛卡尔纤维化}。
\end{definition}

\begin{exercise}\label{L19-2.6}
设 $p\colon\cC\to\cD$ 是拟范畴之间的内纤维化，$f\colon X\to Y$ 是 $p$-笛卡尔态射。对态射 $g\colon X'\to X$，证明 $g$ 是 $p$-笛卡尔的，当且仅当 $f\circ g$ 是。
\end{exercise}

\nitem\label{L19-2.7} 下述结果用提升性质刻画笛卡尔箭头。

\begin{proposition}\label{L19-2.8}（HTT.2.4.1.4, 2.4.1.8, 2.4.4.3）
设 $p\colon\cC\to\cD$ 是拟范畴之间的内纤维化，$f\colon X\to Y$ 是 $\cC$ 中的态射。则下述条件等价：
\begin{enumerate}[label=(\arabic*),leftmargin=2.4em,itemsep=0.1em]
\item 态射 $f$ 是 $p$-笛卡尔的。
\item 当 $n\ge 2$ 且 $\phi|_{\{n-1,n\}}=f$ 时，提升问题
\[
\begin{tikzcd}[column sep=3.8em,row sep=3em]
\Lambda^n_n \arrow[r,"\phi"] \arrow[d] & \cC \arrow[d,"p"]\\
\Delta^n \arrow[r] \arrow[ur,dashed] & \cD
\end{tikzcd}
\]
有解。
\item 当 $n\ge 1$ 且 $\phi|_{\{n\}\times\Delta^1}=f$ 时，提升问题
\[
\begin{tikzcd}[column sep=4.2em,row sep=3.2em]
(\Delta^n\times\{1\})\sqcup_{\partial\Delta^n\times\{1\}}(\partial\Delta^n\times\Delta^1)
 \arrow[r,"\phi"] \arrow[d] & \cC \arrow[d,"p"]\\
\Delta^n\times\Delta^1 \arrow[r] \arrow[ur,dashed] & \cD
\end{tikzcd}
\]
有解。
\end{enumerate}
\end{proposition}

\begin{remark}\label{L19-2.9}
Joyal 与 Lurie 用 (2) 在一般单纯集中定义笛卡尔箭头。
\end{remark}

\nitem\label{L19-2.10} 下述练习说明，笛卡尔箭头在同伦与范畴等价下不变。

\begin{exercise}\label{L19-2.11}
设 $f$ 是 $\cC$ 中的态射，且 $p(f)$ 是同构。证明 $f$ 是 $p$-笛卡尔的，当且仅当它是同构。
\end{exercise}

\begin{exercise}\label{L19-2.12}
若 $f$ 同伦于 $f'$，则 $f$ 是 $p$-笛卡尔的，当且仅当 $f'$ 是。
\end{exercise}

\begin{exercise}\label{L19-2.13}
设给定拟范畴的交换方块
\[
\begin{tikzcd}[column sep=3.4em,row sep=3em]
\cC \arrow[r,"r"] \arrow[d,"p"] & \cC' \arrow[d,"p'"]\\
\cD \arrow[r,"s"] & \cD',
\end{tikzcd}
\]
使得 $r$ 与 $s$ 是等价。则 $\cC$ 中的 $f$ 是 $p$-笛卡尔的，当且仅当 $r(f)$ 是 $p'$-笛卡尔的。
\end{exercise}

\begin{definition}\label{L19-2.14}
设 $p\colon\cC\to\cD$ 是 $\infty$-范畴之间的函子。若对 $p$ 的任意（等价地，某个）拟范畴实现 $p\colon\cC\to\cD$，相应的态射都是 $p$-笛卡尔的，则称 $\cC$ 中的态射 $f$ 是\textbf{$p$-笛卡尔的}。
\end{definition}

\nitem\label{L19-2.15} 事实上，只要限制在范畴纤维化上，笛卡尔纤维化的概念在范畴等价下也不变。为解释这一点，我们需要下述结果，它可由 HTT.2.4.6.5 推出。

\begin{proposition}\label{L19-2.16}
拟范畴之间的笛卡尔纤维化是范畴纤维化。
\end{proposition}

\begin{exercise}\label{L19-2.17}
设给定拟范畴的交换方块
\[
\begin{tikzcd}[column sep=3.4em,row sep=3em]
\cC \arrow[r,"r"] \arrow[d,"p"] & \cC' \arrow[d,"p'"]\\
\cD \arrow[r,"s"] & \cD',
\end{tikzcd}
\]
使得 $r$ 与 $s$ 是等价，且 $p$ 与 $p'$ 是范畴纤维化。证明 $p$ 是笛卡尔纤维化，当且仅当 $p'$ 是。
\end{exercise}

\begin{definition}\label{L19-2.18}
设 $p\colon\cC\to\cD$ 是 $\infty$-范畴之间的函子。若 $p$ 能实现为拟范畴之间的笛卡尔纤维化，则称 $p$ 是\textbf{本质笛卡尔纤维化}。
\end{definition}

\begin{exercise}\label{L19-2.19}
设 $p\colon\cC\to\cD$ 是 $\infty$-范畴之间的函子。证明 $p$ 是本质笛卡尔纤维化，当且仅当对把 $p$ 实现为范畴纤维化 $\cC\to\cD$ 的任意拟范畴实现，相应的函子都是拟范畴之间的笛卡尔纤维化。
\end{exercise}

\nitem\label{L19-2.20} 下述练习解释笛卡尔箭头与局部笛卡尔箭头之间的关系。

\begin{exercise}\label{L19-2.21}
设 $p\colon\cC\to\cD$ 是拟范畴之间的内纤维化，$f\colon X\to Y$ 是 $\cC$ 中的态射。证明 $f$ 是局部 $p$-笛卡尔的，当且仅当 $\cC\times_{\cD}\Delta^1$ 中相应的态射在 $\Delta^1$ 上是笛卡尔的，其中 $\Delta^1\to\cD$ 由 $p(f)$ 给出。
\end{exercise}

\begin{exercise}\label{L19-2.22}
证明任意 $p$-笛卡尔箭头都是局部 $p$-笛卡尔的。
\end{exercise}

\begin{exercise}\label{L19-2.23}
证明 $p$ 是笛卡尔纤维化，当且仅当它是局部笛卡尔纤维化，且局部笛卡尔箭头对复合封闭。
\end{exercise}

\section{笛卡尔箭头的唯一性}

\begin{proposition}\label{L19-3.1}
设 $p\colon\cC\to\cD$ 是拟范畴之间的内纤维化。则 $\cC$ 中的态射 $f\colon X\to Y$ 是局部 $p$-笛卡尔的，当且仅当
\[
\cC_{/Y}\times_{\cD_{/p(Y)}}\{p(f)\}
\]
中相应的对象是终对象。
\end{proposition}

\begin{proof}[证明梗概]
首先可以证明 $\cC_{/Y}\times_{\cD_{/p(Y)}}\{p(f)\}$ 是拟范畴（从而谈论终对象才有意义）。利用练习~\ref{L19-2.21}，可以化归到 $\cD=\Delta^1$、$p(X)=0$、$p(Y)=1$ 的情形。可以识别
\[
\cC_{/f}\to\cC_{/Y}\times_{\cD_{/p(Y)}}\cD_{/p(f)}
\]
为
\[
\bigl(\cC_{/Y}\times_{\cD_{/p(Y)}}\{p(f)\}\bigr)^{/f}\to\cC_{/Y}\times_{\cD_{/p(Y)}}\{p(f)\}.
\]
于是论断可由下述事实推出：$z\in\cE$ 是终对象，当且仅当 $\cE_{/z}\to\cE$ 是平凡 Kan 纤维化（Ker.02HF）。
\end{proof}

\nitem\label{L19-3.2} 特别地，对固定的 $Y\in\cC$ 与 $g\colon x\to p(Y)$，若 $g$ 的（局部）$p$-笛卡尔提升 $X\to Y$ 存在，则它本质唯一。

\section{右纤维化}

\begin{definition}\label{L19-4.1}
若拟范畴之间的函子 $p\colon\cC\to\cD$ 是笛卡尔纤维化，且 $\cC$ 中任意态射都是 $p$-笛卡尔的，则称 $p$ 是\textbf{右纤维化}。若 $p^{\op}$ 是右纤维化，则称 $p$ 是\textbf{左纤维化}。
\end{definition}

\begin{exercise}\label{L19-4.2}
设 $p\colon\cC\to\cD$ 是拟范畴之间的笛卡尔纤维化。证明 $p$ 是右纤维化，当且仅当每个纤维 $\cC_y$（$y\in\cD$）都是 Kan 复形。
\end{exercise}

\begin{exercise}\label{L19-4.3}
设 $p\colon C\to D$ 是单纯集之间的态射，且 $D$ 是拟范畴。证明 $p$ 是右纤维化，当且仅当它对所有 $0<i\le n$ 满足下述右提升性质：
\[
\begin{tikzcd}[column sep=3.8em,row sep=3em]
\Lambda^n_i \arrow[r] \arrow[d] & C \arrow[d,"p"]\\
\Delta^n \arrow[r] \arrow[ur,dashed] & D.
\end{tikzcd}
\]
对偶地，证明 $p$ 是左纤维化，当且仅当它对所有 $0\le i<n$ 满足上述右提升性质。
\end{exercise}

\begin{remark}\label{L19-4.4}
注意当 $D$ 是一般的单纯集时，上述性质仍有意义。事实上，可以用它来定义单纯集之间的右纤维化。
\end{remark}

\begin{exercise}\label{L19-4.5}
证明任意 Kan 纤维化都是右纤维化。
\end{exercise}

\section{例子}

\begin{exercise}\label{L19-5.1}
设 $\cC$ 与 $\cD$ 是拟范畴。证明 $\cC\times\cD\to\cD$ 既是笛卡尔纤维化又是余笛卡尔纤维化，且箭头是（余）笛卡尔的，当且仅当它在 $\cC$ 中的像是同构。
\end{exercise}

\begin{example}\label{L19-5.2}（HTT.2.1.2.2, 4.2.1.6）
设 $u\colon K\to\cC$ 是拟范畴 $\cC$ 中的任意图。则遗忘函子 $\cC_{/u}\to\cC$ 与 $\cC^{/u}\to\cC$ 都是右纤维化。
\end{example}

\begin{example}\label{L19-5.3}（Ker.02VW）
对任意单纯集 $C$，令 $\mathrm{Arr}(C)\coloneqq\Fun([1],C)$ 为 $C$ 中箭头的单纯集。若 $C$ 是拟范畴，则 $\mathrm{Arr}(C)$ 也是。投影
\[
\mathrm{ev}_0\colon\mathrm{Arr}(C)\to C
\]
是笛卡尔纤维化，而投影
\[
\mathrm{ev}_1\colon\mathrm{Arr}(C)\to C
\]
是余笛卡尔纤维化。$\mathrm{Arr}(C)$ 中的态射是 $\mathrm{ev}_0$-笛卡尔的（相应地，$\mathrm{ev}_1$-余笛卡尔的），当且仅当 $\mathrm{ev}_1$（相应地，$\mathrm{ev}_0$）把它送到同构。
\end{example}

\begin{exercise}\label{L19-5.4}
设 $C$ 是拟范畴，$J$、$K$ 是任意单纯集。证明限制函子
\[
p\colon\Fun(J\star K,C)\to\Fun(J,C)
\]
是笛卡尔纤维化，而
\[
q\colon\Fun(J\star K,C)\to\Fun(K,C)
\]
是余笛卡尔纤维化。此外，$\Fun(J\star K,C)$ 中的态射是 $p$-笛卡尔的（相应地，$q$-余笛卡尔的），当且仅当 $q$（相应地，$p$）把它送到同构。提示：使用钝联 $J\Diamond K$ 化归到上例。
\end{exercise}

\begin{example}\label{L19-5.5}（Ker.03JF）
对任意单纯集 $C$，令 $\mathrm{TwArr}(C)$ 为单纯集
\[
\mathrm{TwArr}(C)_n\coloneqq\Hom_{\Set_\Delta}\bigl(\mathrm{N}_{\bullet}([n]^{\op}\star[n]),C\bigr).
\]
若 $C$ 是拟范畴，则 $\mathrm{TwArr}(C)$ 也是，它称为 $C$ 中\textbf{扭箭头}的拟范畴。由 $[n]^{\op}\to[n]^{\op}\star[n]\leftarrow[n]$ 诱导的投影
\[
\mathrm{TwArr}(C)\to C^{\op}\times C
\]
是左纤维化。由此，
\[
\mathrm{TwArr}(C)\to C^{\op},\qquad \mathrm{TwArr}(C)\to C
\]
都是余笛卡尔纤维化。
\end{example}

\begin{exercise}\label{L19-5.6}
$\mathrm{TwArr}(C)\to C^{\op}$ 与 $\mathrm{TwArr}(C)\to C$ 的余笛卡尔箭头分别是什么？
\end{exercise}

\begin{exercise}\label{L19-5.7}
$\mathrm{TwArr}(C^{\op})$ 是什么？
\end{exercise}

\begin{example}\label{L19-5.8}（Ker.01VG）
设 $p\colon\cC\to\cD$ 是拟范畴之间的笛卡尔纤维化，$K$ 是任意单纯集。则
\[
p^{K}\colon\Fun(K,\cC)\to\Fun(K,\cD)
\]
是笛卡尔纤维化，且箭头是 $p^{K}$-笛卡尔的，当且仅当它在任意顶点 $x\in K$ 处的值是 $p$-笛卡尔的。
\end{example}

\begin{exercise}\label{L19-5.9}
设 $\Ring$ 是通常的环范畴，$\LMod$ 是通常的范畴，其对象为对 $(A,M)$，其中 $A$ 是环而 $M$ 是左 $A$-模。$\LMod$ 中的态射 $(A,M)\to(B,N)$ 由环同态 $A\to B$ 与线性映射 $M\to N$ 组成，使得下图交换：
\[
\begin{tikzcd}[column sep=4em,row sep=3em]
A\otimes M \arrow[r] \arrow[d] & M \arrow[d]\\
B\otimes N \arrow[r] & N.
\end{tikzcd}
\]
证明 $\LMod\to\Ring$ 在通常范畴之间既是笛卡尔纤维化又是余笛卡尔纤维化。
\end{exercise}

\appendixsection{局部笛卡尔纤维化与 lax 函子}

设 $\mathsf{S}$ 是 $\infty$-范畴。下述 $(\infty,2)$-范畴之间存在典范等价：
\begin{itemize}[leftmargin=2.2em,itemsep=0.1em]
\item lax 函子 $\mathsf{S}^{\op}\to\Cat_{\infty}$ 的 $(\infty,2)$-范畴，其中 $\Cat_{\infty}$ 是 $(\infty,1)$-范畴的 $(\infty,2)$-范畴；
\item 局部笛卡尔纤维化 $\mathsf{C}\to\mathsf{S}$ 的 $(\infty,2)$-范畴，其中态射是定义在 $\mathsf{S}$ 上、保持 $\mathsf{S}$ 上的笛卡尔箭头的函子。
\end{itemize}

\yizhu{“lax”一词通常译作“松”；本译本沿用第 27 讲的惯例保留英文 lax，以便与“严格”以及“反松”（oplax）等相对术语并列时不易混淆。\textbf{[来源：B]}}

\nitem\label{L19-A.1} \textbf{建议阅读}\quad [Lur09]。

\subsection*{参考文献}
\par\noindent\hangindent=5.2em\hangafter=1 [Lur09]\hspace{0.6em}Jacob Lurie. \emph{$(\infty,2)$-categories and the Goodwillie calculus I}. arXiv preprint arXiv:0905.0462, 2009.
